\section{第5章 定积分}
\subsection{定积分的概念与性质}
\begin{exercise}{557}{基础·积分变量}
设 $I=\int_a^b\arcsin t\dd t$，其中 $-1<a<b<1$。求 $I$ 关于上限 $b$ 和下限 $a$ 的导数。
\end{exercise}
\begin{exercise}{558}{基础·几何意义}
利用半圆面积求 $\int_{-2}^2\sqrt{4-x^2}\dd x$。
\end{exercise}
\begin{exercise}{559}{基础·黎曼和}
将 $\lim_{n\to\infty}\frac1n\sum_{k=1}^n\frac{k}{n}$ 写成定积分并求值。
\end{exercise}
\begin{exercise}{560}{基础·黎曼和}
将 $\sum_{k=1}^n\frac1{n+k}$ 写成黎曼和，求其在 $n\to\infty$ 时的极限。
\end{exercise}
\begin{exercise}{561}{基础·黎曼和}
利用定积分求 $\lim_{n\to\infty}\sum_{k=1}^n\frac{n}{n^2+k^2}$。
\end{exercise}
\begin{exercise}{562}{基础·积分比较}
不求原函数，比较 $\int_0^1x^2\dd x$ 与 $\int_0^1x^3\dd x$ 的大小，并说明依据。
\end{exercise}
\begin{exercise}{563}{基础·积分比较}
不求原函数，比较 $\int_0^1\sqrt{1+x^2}\dd x$ 与 $\int_0^1\sqrt{1+x^3}\dd x$ 的大小。
\end{exercise}
\begin{exercise}{564}{基础·积分估值}
利用 $0\le\cos^2x\le1$，给出 $\int_0^{\frac{\pi}{2}}(1+\cos^2x)\dd x$ 的上、下界。
\end{exercise}
\begin{exercise}{565}{基础·带权积分估值}
利用 $\frac{1}{\sqrt2}\le\frac{1}{\sqrt{1+x}}\le1$（$0\le x\le1$），估计 $\int_0^1\frac{x^2}{\sqrt{1+x}}\dd x$。
\end{exercise}
\begin{exercise}{566}{基础·三角积分估值}
利用 $\frac{\sqrt2}{2}\le\sin x\le1$，估计 $\int_{\frac{\pi}{4}}^{\frac{\pi}{2}}\sin x\dd x$。
\end{exercise}
\begin{exercise}{567}{常规·积分中值定理}
设 $f$ 在 $[0,1]$ 上连续。用积分中值定理说明存在 $\xi\in[0,1]$ 使 $\int_0^1f(x)\dd x=f(\xi)$；当 $f(x)=x^2$ 时求 $\xi$。
\end{exercise}
\begin{exercise}{568}{常规·积分平均值}
设 $f$ 在 $[0,1]$ 上连续且单调减少。证明 $\int_0^xf(t)\dd t\ge xf(x)$（$0<x\le1$）。提示：比较 $0\le t\le x$ 时的 $f(t)$ 与 $f(x)$。
\end{exercise}
\begin{exercise}{569}{基础·函数平均值}
求 $f(x)=\sqrt{4-x^2}$ 在 $[-2,2]$ 上的平均值。
\end{exercise}
\begin{exercise}{570}{基础·黎曼和变式}
把 $\lim_{n\to\infty}\sum_{k=1}^n\frac{k}{n^2+k^2}$ 写成定积分，并用换元法求值。
\end{exercise}
\begin{exercise}{571}{常规·积分不等式}
设 $f$ 在 $[0,1]$ 上连续且 $0\le f\le1$。证明 $0\le\int_0^1f^2(x)\dd x\le\int_0^1f(x)\dd x$。提示：先比较被积函数。
\end{exercise}
\subsection{微积分基本公式}
\begin{exercise}{572}{基础·双变限求导}
求 $y=\int_x^{x^2}\sin(t^2)\dd t$ 的导数。
\end{exercise}
\begin{exercise}{573}{基础·积分上限求导}
设 $y=\int_0^{x^2}\sqrt{1+t^2}\dd t$，求 $y'$ 及 $y'(1)$。
\end{exercise}
\begin{exercise}{574}{基础·双变限求导}
求 $y=\int_x^{2x}\frac{\dd t}{1+t^2}$ 的导数及 $y'(0)$。
\end{exercise}
\begin{exercise}{575}{基础·变限积分极限}
求 $\lim_{x\to0}\frac{\int_0^x\cos(t^2)\dd t}{x}$。
\end{exercise}
\begin{exercise}{576}{基础·变限积分极限}
利用洛必达法则求 $\lim_{x\to0}\frac{\int_0^x\arctan t\dd t}{x^2}$。
\end{exercise}
\begin{exercise}{577}{常规·积分商极限}
求 $\lim_{x\to0}\frac{(\int_0^xe^{t^2}\dd t)^2}{x^2}$。提示：先求积分与 $x$ 的比值极限。
\end{exercise}
\begin{exercise}{578}{常规·含参变限积分}
设 $f$ 在原点附近连续，求 $\lim_{x\to0}\frac{\int_0^x tf(t)\dd t}{x^2}$，结果用 $f(0)$ 表示。
\end{exercise}
\begin{exercise}{579}{基础·无穷远积分极限}
利用洛必达法则求 $\lim_{x\to+\infty}\frac{\int_0^x\arctan t\dd t}{x}$。
\end{exercise}
\begin{exercise}{580}{常规·隐式变限积分}
设 $\int_0^ye^t\dd t+\int_0^x\cos t\dd t=0$ 在原点附近确定 $y=y(x)$。求 $y'(0)$。
\end{exercise}
\begin{exercise}{581}{基础·变限函数单调性}
设 $f(x)=\int_0^xe^{-\frac{t^2}{2}}\dd t$，求 $f'(x)$ 并判断 $f$ 的单调性。
\end{exercise}
\begin{exercise}{582}{常规·求导恢复数据}
设 $f$ 连续，$\varphi(x)=x\int_0^{x^2}f(t)\dd t$，已知 $\varphi(1)=1$、$\varphi'(1)=5$。求 $f(1)$。
\end{exercise}
\begin{exercise}{583}{常规·积分函数零点}
设 $F(x)=\int_0^x(1+t^2)\dd t-1$。证明 $F$ 在 $(0,1)$ 内恰有一个零点。提示：计算端点值和导数。
\end{exercise}
\begin{exercise}{584}{基础·牛顿莱布尼茨公式}
计算 $\int_{-1}^{\frac{1}{2}}e^{2x}\dd x$。
\end{exercise}
\begin{exercise}{585}{基础·基本公式}
计算 $\int_0^2\frac{x}{\sqrt{1+x^2}}\dd x$。
\end{exercise}
\begin{exercise}{586}{基础·部分分式积分}
计算 $\int_0^1\frac{\dd x}{x^2+5x+6}$。
\end{exercise}
\begin{exercise}{587}{基础·绝对值积分}
计算 $\int_{-1}^3|2-x|\dd x$，先确定绝对值内表达式改变符号的位置。
\end{exercise}
\begin{exercise}{588}{基础·三角绝对值积分}
计算 $\int_0^{2\pi}|\sin x|\dd x$。
\end{exercise}
\begin{exercise}{589}{常规·分段变限积分}
设 $f(x)=2x$（$0\le x\le1$），$f(x)=3$（$1<x\le2$）。求 $P(x)=\int_0^xf(t)\dd t$ 的分段表达式。
\end{exercise}
\begin{exercise}{590}{常规·积分常数恢复}
连续函数 $f$ 在 $[0,1]$ 上满足 $f(x)=x+\frac12\int_0^1f(t)\dd t$。令 $A=\int_0^1f(t)\dd t$，求 $A$ 和 $f(x)$。
\end{exercise}
\begin{exercise}{591}{常规·积分方程}
连续函数 $f$ 在 $[0,1]$ 上满足 $f(x)=x\int_0^1f(t)\dd t-1$。求 $\int_0^1f(t)\dd t$ 和 $f(x)$。
\end{exercise}
\begin{exercise}{592}{基础·复合换元}
计算 $\int_1^2(2x+1)(x^2+x)\dd x$，提示：令 $u=x^2+x$。
\end{exercise}
\begin{exercise}{593}{基础·变限函数极值}
设 $F(x)=\int_0^x(t-1)e^{-t}\dd t$，$x\ge0$。求 $F$ 的单调区间及极小值。
\end{exercise}
\begin{exercise}{594}{常规·变限积分小综合}
设 $F(x)=\int_0^{x^2}(1+t)\dd t$。求 $F'(x)$，并计算 $\int_0^1F'(x)\dd x$。
\end{exercise}
\subsection{定积分的换元积分法和分部积分法}
\begin{exercise}{595}{基础·根式换元}
计算 $\int_4^9\frac{\dd x}{\sqrt x-1}$，提示：令 $u=\sqrt x$。
\end{exercise}
\begin{exercise}{596}{基础·根式换元}
计算 $\int_0^1x\sqrt{3-2x}\dd x$。
\end{exercise}
\begin{exercise}{597}{基础·对数换元}
计算 $\int_1^e\frac{\dd x}{x\sqrt{1+\ln x}}$。
\end{exercise}
\begin{exercise}{598}{基础·指数换元}
计算 $\int_0^1\frac{e^x}{\sqrt{1+e^x}}\dd x$。
\end{exercise}
\begin{exercise}{599}{基础·三角换元}
计算 $\int_0^{\frac{\sqrt2}{3}}\sqrt{2-9x^2}\dd x$。
\end{exercise}
\begin{exercise}{600}{基础·三角换元}
计算 $\int_0^1\frac{\dd x}{\sqrt{1-x^2}}$，按上端点极限理解积分。
\end{exercise}
\begin{exercise}{601}{基础·根式换元}
计算 $\int_1^2x\sqrt{x^2-1}\dd x$。
\end{exercise}
\begin{exercise}{602}{常规·根式与对称性}
利用偶函数性质及三角换元，计算 $\int_{-1}^1\frac{\dd x}{1+\sqrt{1-x^2}}$。
\end{exercise}
\begin{exercise}{603}{基础·配方与几何意义}
将 $2x-x^2$ 配方，计算 $\int_0^1\sqrt{2x-x^2}\dd x$。
\end{exercise}
\begin{exercise}{604}{基础·指数换元}
计算 $\int_0^{\ln2}\frac{e^{2x}}{1+e^x}\dd x$，提示：令 $u=e^x$。
\end{exercise}
\begin{exercise}{605}{基础·平移分段积分}
设 $f(t)=1+t^2$（$t<0$），$f(t)=e^{-t}$（$t\ge0$）。计算 $\int_1^3f(x-2)\dd x$。
\end{exercise}
\begin{exercise}{606}{基础·平移与奇偶性}
设 $f(t)=te^{t^2}$（$-1\le t\le1$）。求 $\int_0^2f(x-1)\dd x$。
\end{exercise}
\begin{exercise}{607}{基础·积分函数恢复}
设 $f$ 连续且 $\int_0^xf(t)\dd t=e^{-2x}-1$。求 $f(x)$ 及 $\int_0^1f(x)\dd x$。
\end{exercise}
\begin{exercise}{608}{常规·二次积分求导}
设 $f$ 连续且 $\int_0^x(x-t)f(t)\dd t=e^x-x-1$。对等式求导两次，求 $f(x)$。
\end{exercise}
\begin{exercise}{609}{基础·奇函数积分}
计算 $\int_{-\frac{\pi}{2}}^{\frac{\pi}{2}}\frac{x}{1+\cos x}\dd x$。
\end{exercise}
\begin{exercise}{610}{基础·对数奇偶性}
利用 $(\sqrt{1+x^2}+x)(\sqrt{1+x^2}-x)=1$，计算 $\int_{-1}^1\ln(x+\sqrt{1+x^2})\dd x$。
\end{exercise}
\begin{exercise}{611}{常规·对称换元}
用 $x\mapsto-x$ 计算 $\int_{-1}^1\frac{\dd x}{1+e^x}$。
\end{exercise}
\begin{exercise}{612}{常规·反三角对称性}
利用 $\arccos(-x)=\pi-\arccos x$，计算 $\int_{-1}^1\cos x\arccos x\dd x$。
\end{exercise}
\begin{exercise}{613}{基础·奇偶分解}
计算 $\int_{-1}^1(x^4+\sin x)\dd x$。
\end{exercise}
\begin{exercise}{614}{基础·周期与绝对值}
计算 $\int_0^{2\pi}\sqrt{1-\sin^2x}\dd x$，注意根式等于 $|\cos x|$。
\end{exercise}
\begin{exercise}{615}{基础·三角积分}
计算 $\int_0^{\frac{\pi}{4}}\sin x\cos x\dd x$。
\end{exercise}
\begin{exercise}{616}{常规·对称换元}
利用 $x\mapsto\frac{\pi}{2}-x$，计算 $\int_0^{\frac{\pi}{2}}\frac{\sin x}{\sin x+\cos x}\dd x$。
\end{exercise}
\begin{exercise}{617}{基础·分部积分}
计算 $\int_0^1\ln(1+x^2)\dd x$。
\end{exercise}
\begin{exercise}{618}{基础·反三角分部积分}
计算 $\int_0^{\frac{1}{2}}\arcsin x\dd x$。
\end{exercise}
\begin{exercise}{619}{常规·反三角分部积分}
用分部积分法计算 $\int_{-1}^1x\arccos x\dd x$。
\end{exercise}
\begin{exercise}{620}{基础·对数分部积分}
计算 $\int_0^1x\ln(1+x)\dd x$。
\end{exercise}
\begin{exercise}{621}{常规·循环分部积分}
计算 $\int_0^{\frac{\pi}{2}}e^{-x}\sin2x\dd x$。
\end{exercise}
\begin{exercise}{622}{常规·换元与分部积分}
计算 $\int_0^1e^{-\sqrt x}\dd x$，先令 $t=\sqrt x$。
\end{exercise}
\begin{exercise}{623}{基础·三角换元}
计算 $\int_0^1\frac{\arcsin x}{\sqrt{1-x^2}}\dd x$，按上端点极限理解积分。
\end{exercise}
\begin{exercise}{624}{常规·反三角分部积分}
计算 $\int_0^1\arctan x\dd x$。
\end{exercise}
\begin{exercise}{625}{基础·倒数换元}
计算 $\int_1^2x^{-2}e^{\frac{1}{x}}\dd x$。
\end{exercise}
\begin{exercise}{626}{常规·根式换元与分部积分}
计算 $\int_0^{\pi^2}\cos\sqrt x\dd x$，先令 $t=\sqrt x$。
\end{exercise}
\begin{exercise}{627}{基础·端点数据积分}
设 $f\in C^2[0,2]$，$f(0)=1$、$f(2)=3$、$f'(2)=5$。求 $\int_0^2xf''(x)\dd x$。
\end{exercise}
\begin{exercise}{628}{基础·周期变限积分}
设 $F(x)=\int_x^{x+2\pi}\sin t\dd t$。求 $F'(x)$ 及 $F(x)$。
\end{exercise}
\begin{exercise}{629}{常规·已知原函数}
已知 $F(x)=e^{x^2}$ 是 $f(x)$ 的一个原函数。计算 $\int_0^1xf'(x)\dd x$，提示：先分部积分。
\end{exercise}
\begin{exercise}{630}{常规·积分常数恢复}
设连续函数 $f$ 满足 $f(x)=x+\int_0^1tf(t)\dd t$。求 $f(x)$。提示：把积分记为常数后代回求值。
\end{exercise}
\begin{exercise}{631}{常规·变限积分分部计算}
设 $f(x)=\int_0^x\sin t\dd t$。计算 $\int_0^\pi f(x)\dd x$。
\end{exercise}
\begin{exercise}{632}{常规·变限函数分部积分}
设 $f(x)=\int_0^xe^{-t^2}\dd t$。计算 $\int_0^1xf(x)\dd x$，结果允许含 $\int_0^1e^{-t^2}\dd t$。
\end{exercise}
\begin{exercise}{633}{常规·求导与基本公式}
设 $f(x)=\int_0^x(1+t^2)\dd t$。计算 $\int_0^1xf'(x)\dd x$，并用分部积分核验。
\end{exercise}
\begin{exercise}{634}{基础·对称换元变式}
用 $x\mapsto\frac{\pi}{2}-x$ 计算 $\int_0^{\frac{\pi}{2}}\frac{\cos x}{\sin x+\cos x}\dd x$。
\end{exercise}
\begin{exercise}{635}{常规·分部积分小综合}
设 $I_n=\int_0^1x^ne^x\dd x$，$n\ge1$ 为整数。用分部积分得到 $I_n=e-nI_{n-1}$，并据此求 $I_2$。
\end{exercise}
\subsection{反常积分}
\begin{exercise}{636}{基础·无穷积分判敛}
分别判断 $\int_1^{+\infty}x^{-\frac{1}{2}}\dd x$ 与 $\int_1^{+\infty}x^{-2}\dd x$ 的敛散性，对收敛者求值。
\end{exercise}
\begin{exercise}{637}{基础·对数反常积分}
计算 $\int_2^{+\infty}\frac{\dd x}{x\ln^2x}$。
\end{exercise}
\begin{exercise}{638}{基础·全轴积分}
计算 $\int_{-\infty}^{+\infty}\frac{\dd x}{x^2+4x+5}$，先对分母配方。
\end{exercise}
\begin{exercise}{639}{基础·部分分式反常积分}
计算 $\int_1^{+\infty}\frac{\dd x}{x^2+x}$，先在 $[1,R]$ 上积分再取极限。
\end{exercise}
\begin{exercise}{640}{基础·根式反常积分}
计算 $\int_1^{+\infty}x^{-\frac{3}{2}}\dd x$。
\end{exercise}
\begin{exercise}{641}{常规·全轴换元积分}
令 $u=e^x$，计算 $\int_{-\infty}^{+\infty}\frac{e^x}{1+e^{2x}}\dd x$。
\end{exercise}
\begin{exercise}{642}{基础·指数反常积分}
计算 $\int_0^{+\infty}e^{-2x}\dd x$。
\end{exercise}
\begin{exercise}{643}{常规·反常分部积分}
计算 $\int_1^{+\infty}\frac{\ln x}{x^2}\dd x$，写出分部积分中无穷端边界项的极限。
\end{exercise}
\begin{exercise}{644}{常规·三角代换}
计算 $\int_0^{+\infty}\frac{\dd x}{(1+x^2)^2}$，提示：令 $x=\tan t$。
\end{exercise}
\begin{exercise}{645}{基础·反三角反常积分}
计算 $\int_1^{+\infty}\frac{\dd x}{1+x^2}$。
\end{exercise}
\begin{exercise}{646}{常规·反常分部积分}
计算 $\int_0^{+\infty}xe^{-x}\dd x$。
\end{exercise}
\begin{exercise}{647}{基础·全轴绝对收敛积分}
判断 $\int_{-\infty}^{+\infty}xe^{-x^2}\dd x$ 是否绝对收敛，并求其值。提示：分别考察两个半轴。
\end{exercise}
\begin{exercise}{648}{基础·比较判别法}
利用 $\ln x<x$（$x\ge2$），判断 $\int_2^{+\infty}\frac{\dd x}{\ln x}$ 的敛散性。
\end{exercise}
\begin{exercise}{649}{常规·高斯型积分判敛}
用 $x\ge1$ 时 $e^{-x^2}\le e^{-x}$，证明 $\int_0^{+\infty}e^{-x^2}\dd x$ 收敛，无需求准确值。
\end{exercise}
\begin{exercise}{650}{基础·指数反常积分}
计算 $\int_{-\infty}^0e^{2t}\dd t$。
\end{exercise}
\begin{exercise}{651}{基础·瑕积分判敛}
判断 $\int_0^1x^{-\frac{1}{2}}\dd x$ 与 $\int_0^1x^{-2}\dd x$ 的敛散性，对收敛者求值。
\end{exercise}
\begin{exercise}{652}{常规·对数瑕积分}
计算 $\int_0^1x\ln x\dd x$，写出 $x\to0^+$ 时边界项的极限。
\end{exercise}
\begin{exercise}{653}{基础·根式瑕积分}
计算 $\int_0^2\frac{\dd x}{\sqrt{4-x^2}}$，用极限处理上端点。
\end{exercise}
\begin{exercise}{654}{基础·瑕点识别}
判断 $\int_0^1\frac{\dd x}{x(\ln x)^2}$ 的敛散性，分别考察 $0$ 和 $1$ 两个端点。
\end{exercise}
\begin{exercise}{655}{基础·幂次判别}
讨论 $\int_0^1x^{-p}\dd x$ 随实参数 $p$ 的敛散性，收敛时求值。
\end{exercise}
\begin{exercise}{656}{常规·倒数换元}
计算 $\int_1^{+\infty}\frac{\ln x}{x^2}\dd x$，这次使用 $t=\frac{1}{x}$ 将其化为瑕积分。
\end{exercise}
\begin{exercise}{657}{基础·分段反常积分}
设 $f(x)=x^{-\frac{1}{2}}$（$0<x\le1$），$f(x)=x^{-2}$（$x>1$）。计算 $\int_0^{+\infty}f(x)\dd x$。
\end{exercise}
\begin{exercise}{658}{常规·内部瑕点}
判断 $\int_{-1}^1x^{-2}\dd x$ 是否存在，先在 $0$ 处分成两个反常积分。
\end{exercise}
\begin{exercise}{659}{基础·双端根式瑕积分}
令 $x=\sin^2t$，计算 $\int_0^1\frac{\dd x}{\sqrt{x(1-x)}}$。
\end{exercise}
\begin{exercise}{660}{常规·对数换元小综合}
判断 $\int_e^{+\infty}\frac{\dd x}{x\ln x}$ 与 $\int_e^{+\infty}\frac{\dd x}{x(\ln x)^2}$ 的敛散性，对收敛者求值。
\end{exercise}
